% ============================================================================
\clearpage
\section*{v4.10.0-r50 Source-Only: A Medium from Light-Like Streams (H120) and a Vacuum Bound on the Directed Coupling}
\addcontentsline{toc}{section}{v4.10.0-r50 Source-Only: Light-Like Medium and Vacuum Bound (H120)}
\label{sec:v4100-r50-null-medium}
% ============================================================================

\paragraph{Purpose.}
By v4.10.0-r47 and r49 the OP5 comparison model needs a medium that slows the Higgs radius and both mediators in every direction.  The author
asked whether such a medium could be an object of the photon frame P that appears bent in the other frames, and chose to start with ordinary
boosts (reading a).  Since a boost is only a change of description, the bending is made measurable through invariants \(u\cdot l\) between an
observer and the light-like vectors \(l\) of the medium.  This layer records hypothesis H120, fixed before the computation.  H120a is not blind;
H120b and H120c are blind in their values.  Earlier layers are not rewritten.

\noindent\textbf{Source-only status.}
This v4.10.0-r50 layer is source-only.  It does not build a PDF, does not assign a new DOI,
and does not authorize a public release.  The v5.3.0 Version DOI
\href{https://doi.org/10.5281/zenodo.24434403}{10.5281/zenodo.24434403} and the Concept DOI
\href{https://doi.org/10.5281/zenodo.19210728}{10.5281/zenodo.19210728} are unchanged by it.

\subsection*{v4.10.0-r50 setting}
\label{subsec:v4100-r50-setting}

All three fields propagate on one effective metric \(G^{\mu\nu}=\eta^{\mu\nu}+\kappa\sum_sw_s\,l_s^\mu l_s^\nu\) with light-like \(l_s\) and \(\kappa>0\): one stream \(l=(1,0,0,1)\), or two
counter-streams \(l_\pm=(1,0,0,\pm1)\) with weights \(a,b\).  An observer with rapidity \(\eta\) along the stream axis sees the weights \(a\,e^{-2\eta}\) and \(b\,e^{2\eta}\); the
photon frame is the limit \(\eta\to\infty\).  The cone limit is evaluated as in r47 with \(c_h^2=r_0=r_m=r(\theta)\), the direction-dependent squared speed.

\subsection*{v4.10.0-r50 results}
\label{subsec:v4100-r50-results}

\paragraph{Theoremlet v4.10.0-r50.1 (one stream, H120a confirmed).}
For a single light-like stream \(v=1\) is a root in the direction of \(l\) for every \(\kappa\), and there the largest squared characteristic speed is \(1.1995\), for \(\kappa w\) from
\(0.1\) to \(100\).  A medium made of one stream of the photon frame cannot carry the slowing.
\emph{Proof.} Substitution in the dispersion relation; \(j_m^2\ge\tfrac12g_m^2>0\) for all \(\alpha\); numerical eigenvalues (script \texttt{k240}).  \(\square\)

\paragraph{Theoremlet v4.10.0-r50.2 (two counter-streams, H120b partially confirmed).}
Two counter-streams keep the cone limit in all directions for 35 of 49 parameter sets on a conservative grid, for example \(0.803\) at \(\kappa a=\kappa b=0.3\).  For
an observer moving along the stream axis the largest slowing \(\max_\theta(1-r)\) is \(0.75\), \(0.52\), \(0.69\), \(0.94\) at \(\eta=0,\,0.6,\,1,\,2\) and tends to one, while the
cone limit gives \(0.80\), \(0.93\), \(1.05\), \(1.18\) and tends to \(1.1995\).  The medium does not vanish towards the photon frame; it degenerates into one stream.
\emph{Proof.} Numerical roots and eigenvalues (script \texttt{k240}).  \(\square\)

\paragraph{Theoremlet v4.10.0-r50.3 (minimal slowing, H120c confirmed).}
For a universal slowing \(r=1-\varepsilon\) the upper cone condition (times \(1-r\)) is exactly \(\varepsilon^2-g_m^2\bigl(\tfrac23-\tfrac23\alpha+\tfrac76\alpha^2\bigr)\ge0\) (with \(Z=1\), \(g_0=g_m\)), sharpest at
\(\alpha=1\).  Hence \(\varepsilon\ge\sqrt{7g_m^2/6}\), confirmed numerically at \(g_m^2=\tfrac3{100},10^{-4},10^{-8}\).
\emph{Proof.} Symbolic expansion (script \texttt{k240}).  \(\square\)

\paragraph{Remark v4.10.0-r50.4 (reading).}
In reading (a) the idea takes an inverted form: structures of the photon frame, two light-like counter-streams, form for observers near their
balance a genuine medium with a rest frame (where the weights are equal, as in r41), whereas in the photon frame itself they are a single stream.
Thesis T9 (\(\Lambda=0\) in P) is not supported by it.  The bound r50.3 holds for every medium with a universal slowing in vacuum: with a laboratory bound
\(\varepsilon\le10^{-13}\) on Lorentz violation (H78) it gives \(g_m^2\le\tfrac67\cdot10^{-26}\approx8.6\times10^{-27}\), so the comparison model with \(g_m^2=\tfrac3{100}\) is excluded in vacuum by about
25 orders of magnitude, and an observable directed contrast needs a material whose rest frame is the preparation.  The bound \(10^{-13}\) concerns
fermions and is adopted here for the Higgs sector without separate evidence.  Found afterwards and not fixed beforehand: at \(\eta=0.6\), close to the
rapidity of M (\(\operatorname{artanh}\beta_{\rm sym}=0.618\)), the limit still holds; the rapidity here runs along the stream axis, that of OP1 across the axis of R,
so this is no confirmation.

\subsection*{v4.10.0-r50 frames}
\label{subsec:v4100-r50-frames}

Speeds refer to the flat metric of the matter frame M with \(c=1\).  The photon frame P enters as the limit \(\eta\to\infty\) of ordinary boosts (reading a); the
frame map \(J_B\) (reading b) is not used.  Time reversal and the tachyon map are not used.

\subsection*{v4.10.0-r50 register entries}
\label{subsec:v4100-r50-register}

\begin{center}
\small
\begin{tabular}{@{}>{\raggedright\arraybackslash}p{0.34\textwidth}>{\raggedright\arraybackslash}p{0.58\textwidth}@{}}
\toprule
statement & status after v4.10.0-r50\\
\midrule
H120a (one light-like stream) & confirmed: cannot carry the slowing\\
H120b (two counter-streams) & partially confirmed: cone limit holds near balance; not effective only in P\\
H120c (vacuum bound) & confirmed: \(\varepsilon\ge\sqrt{7g_m^2/6}\); with \(\varepsilon\le10^{-13}\), \(g_m^2\lesssim8.6\times10^{-27}\)\\
OP5 & open\\
\bottomrule
\end{tabular}
\end{center}

\subsection*{v4.10.0-r50 notes}
\label{subsec:v4100-r50-notes}

\texttt{llm/\allowbreak Vorab\_\allowbreak Notiz\_\allowbreak H120\_\allowbreak Nullmedium.md} (commit d3810e58), \texttt{llm/\allowbreak Ergebnis\_\allowbreak H120\_\allowbreak Nullmedium.md},
script \texttt{llm/\allowbreak k240\_\allowbreak h120\_\allowbreak null\_\allowbreak medium.py}, output \texttt{llm/\allowbreak k240\_\allowbreak out.txt}.

\subsection*{v4.10.0-r50 ledger and guard update}
\label{subsec:v4100-r50-ledger-guard}

\[
R_{\rm chk}^{v4100\text{-}r50}=(1,1,0)
\]
for (status recorded; computations checked in the cited script; a material or vacuum medium of the required kind not identified).

The layer spends no guard:
\[
R_{\rm glob}^{(v4100\text{-}r50,+)}=(0,1,1,1,1),
\qquad
\Gamma_{\rm guard}^{v4100\text{-}r50}=(1,1,1,1,1),
\qquad
\beta_4=1/256\ \text{unspent}.
\]
The layer does not imply exact finite \(C5\) zero closure, exact all-readout terminality, \(C6/T2\), a global photon-clock/RH theorem,
or global all-readout closure.  OP1--OP5, C5 and C6/T2 remain open.
